\section{Terminal trajectories and learned execution policies}\label{sec:C23:terminal}
Terminal agents expose a broad action space through a compact interface. The same shell can inspect files, run tests, modify dependencies, and invoke network tools. This flexibility reduces the number of specialized APIs needed, but it makes semantic validation harder. A syntactically valid command can have effects far beyond the intended task. The sandbox and permission layer must mediate those effects independently of the model's stated intent.

Long terminal trajectories contain more than command examples. They include evidence gathering, hypothesis revision, failed attempts, recovery, and final artifact validation. TerminalHorizon explicitly studies synthesis of task environments that preserve the work needed for these behaviors~\cite{corpus261013166}. Its environment-reconstruction mechanism is relevant because a synthetic setup can accidentally solve the difficult part before the agent begins. A training example then rewards superficial completion rather than the intended exploration.

The current v2 separates a task package into public instruction and environment and private reference solution, verifier and construction configuration. Its independent challenger is uninvolved in construction and solves from the public interface; trace-based diagnosis then guides reconstruction of work accidentally simplified during setup. Progressive expansion adds requirements under the same goal that depend on earlier results and feedback~\cite{corpus261013166}. This is more specific than increasing a maximum turn count. Challenger success alone does not certify absence of every shortcut, and generated verifier agreement remains conditional on its task specification. The appropriate ablation changes reconstruction or expansion separately while preserving model, training tokens and evaluation access, and reports both task coverage and the work demonstrably left for the learner.

The causal comparison should distinguish richer trajectories from easier environments. Keep initial information, required work, and evaluator access comparable across data-generation variants. Audit whether task descriptions leak a solution, whether setup scripts precompute key artifacts, and whether generated tests reveal hidden requirements. Supervised fine-tuning gains on multiple benchmarks can be informative while leaving these construction confounds unresolved.

Trajectory learning also depends on actual behavior conditioning. If a training pipeline reconstructs context after the fact and removes information the policy originally saw, logged actions no longer correspond to the assumed policy inputs. MemoryWalker addresses this class of mismatch~\cite{j2026memorywalker}. For coding agents, record context edits, tool-output truncation, and environment version along with actions. A repository checkout and final patch are insufficient to reconstruct the decision process.

Evaluation should report progress through the full task rather than command validity alone. Measure artifact success, constraints, execution cost, recovery, and whether the final output identifies tested and untested behavior. A terminal agent that produces a plausible explanation after failing to run tests should not receive the same completion status as one with verified results. Conversely, an infrastructure outage should remain visible as an environmental failure rather than be silently removed from the denominator.


The inspected mini-swe-agent release 2.4.6, dated 23 July 2026, explicitly counts billed calls that fail parsing against the cost limit\cite{executionrelease08}. This small implementation detail materially changes budget-matched evaluation: excluding paid malformed responses would give a system additional effective resources while reporting the same nominal limit. It is a concrete example of why a minimal loop remains useful as a research baseline but still requires release-level inspection. The relevant experiment compares identical malformed-output schedules before and after the accounting change, measuring budget overshoot and completed work rather than inferring a reasoning improvement from the release note.

The pinned implementation confirms a further boundary. \path{DefaultAgent.query} checks already accumulated cost before the next model call; successful calls add returned cost, while the \path{FormatError} recovery branch adds reported billed cost from the exception record~\cite{cumb2026minisweloop}. This is post-charge accounting, not the reservation protocol of Chapter~\ref{ch:C20}. For one serial query loop, assume admission and cost updates are serialized, at most one charge remains unrecorded in flight, and each charge is recorded before admitting the next call. If $B>0$ is the active threshold, the initial total is below $B$, and each newly admitted call costs at most $c_{\max}$, the pre-call check bounds the charged total by $B+c_{\max}$, rather than by $B$. Parallel callers can all pass a stale total, so the one-call bound is not a global concurrency guarantee. Without a known call-cost ceiling there is no finite overshoot bound from that check alone. Tool fees and environment charges outside this counter require separate accounting. A reproducible budget comparison must therefore distinguish a declared stopping threshold from a guaranteed maximum expenditure.
